\section{Introduction}
\label{sec:intro}

Serving a large language model (LLM) is a queueing problem with an unusual resource: the
key--value (KV) cache, which grows with every token a request generates and must stay in GPU
memory until the request finishes. Modern serving engines---\vllm~\cite{kwon2023vllm},
\sglang~\cite{zheng2024sglang}, and \trt~\cite{trtllm}---batch requests continuously at
iteration granularity~\cite{yu2022orca} and manage the KV cache in blocks, which lets them
keep hundreds of requests in flight on one GPU. When the KV cache or the scheduler's slots
run out, however, each engine must choose: admit a new request and risk evicting a running
one, or hold the new request back. These choices are made inside the engine, are rarely
documented as part of the engine's contract, and are not exposed through the OpenAI-compatible
APIs that applications use.

Practitioners choosing an engine typically rely on published per-token throughput figures or
on vendor benchmarks run at a fixed, comfortable load. Such figures answer the question
``how fast can this engine generate tokens?'' but not the question operators actually face:
``what will my users experience when traffic is heavy, request lengths vary, and memory is
tight?'' The answer depends on several metrics at once---how long a user waits for the first
token, whether the stream stalls midway, how long the whole answer takes, how many requests
fail---and on which of these the application's service-level objective (SLO) is bound by.

This paper compares \vllm, \sglang, and \trt empirically under conditions that expose their
scheduling behaviour. We drive all three through their own OpenAI-compatible serving
frontends with a common open-loop harness, \sysname, using realistic workloads assembled from
public corpora, Poisson and bursty arrival processes, and three models whose per-token KV
footprint differs by up to $6.25\times$, including one model with full multi-head attention
(MHA) and two with grouped-query attention (GQA)~\cite{ainslie2023gqa}. We organise the study
around four research questions:

\begin{description}
  \item[\rqref{1}] \emph{Engine behaviour under load.} How do the engines compare in
  responsiveness, streaming quality, end-to-end latency, throughput, reliability, and SLO
  goodput as load and KV-cache pressure increase, and what scheduler behaviour explains the
  differences?
  \item[\rqref{2}] \emph{Prefill handling.} How do prefix caching and chunked prefill change
  latency under load on each engine?
  \item[\rqref{3}] \emph{Speculative decoding.} Does n-gram speculative decoding improve
  latency and throughput, and is the effect consistent across engines?
  \item[\rqref{4}] \emph{Hardware generalisation.} Do the \rqref{1} findings hold on a
  different hardware platform with a much larger, unified memory?
\end{description}

\paragraph{Findings.}
\todo{finalise against \cref{sec:results}; the statements below reflect the current analysis}
\begin{itemize}
  \item At light load the engines are interchangeable: decode speed differs by at most 6\%.
  They diverge once a resource saturates, and the dominant difference is \emph{where} each
  engine makes requests wait: \vllm admits early and later stalls running streams,
  \sglang holds requests before their first token and then streams smoothly, and \trt admits
  the fewest requests.
  \item The engine with the fastest token generation (\trt) has the lowest SLO goodput on
  heterogeneous traffic, yet the lowest end-to-end latency on uniform multi-turn traffic.
  Under the SLOs we study, goodput is decided by time to first token, not by token speed.
  \item Raising the scheduler's sequence limit doubles goodput when slots are the bottleneck
  and has no effect when the KV cache or the engine's admission policy is.
  \item Performance features are engine-specific: prefix caching helps all engines alike,
  whereas chunked prefill and n-gram speculative decoding help one engine and hurt another.
  \item The engine ranking under KV pressure is preserved on a second platform with
  $2.7\times$ the memory and newer engine versions.
\end{itemize}

\paragraph{Contributions.}
\begin{enumerate}
  \item A controlled, multi-metric comparison of three production LLM serving engines under
  load, covering latency, streaming quality, throughput, reliability, and SLO goodput
  (\cref{sec:results}).
  \item An explanation of the observed differences in terms of each engine's admission and
  eviction behaviour, the scheduler's sequence limit, and the KV-cache footprint of the model.
  \item A set of methodology lessons for benchmarking LLM serving, each of which changes a
  conclusion if ignored (\cref{sec:discussion}).
  \item A decision guide for practitioners, and an open benchmark harness with all
  per-request data (\cref{sec:conclusion}).
\end{enumerate}

The remainder of the paper is organised as follows. \Cref{sec:background} summarises the
serving mechanisms that matter for our results. \Cref{sec:method} describes the harness,
engines, models, workloads, and metrics. \Cref{sec:results} answers the research questions.
\Cref{sec:discussion} discusses implications for practitioners and for benchmarking,
\cref{sec:threats} threats to validity, and \cref{sec:related} related work.
